\section{Whittaker 系数}\srcpage{40}

谱作用在各种设定下给出 $D_{\qc}(\LocSys_{\Gd})$ 在 $D(\Bun_G)$ 上的一个作用。在理想的情形下，这会使 $D(\Bun_G)$ 成为秩 $1$ 的自由模。尽管这一点会略有失效，仍然有必要造出所预期的基元素。换言之，我们寻找 $D(\Bun_G)$ 的一个特殊对象。这就是 Whittaker 层（亦称 Poincaré 对象），其定义如下。

回忆我们固定了 $\OmX$ 的一个平方根 $\Theta\in\Pic(X)$。特别地，把 Borel 子群 $B$ 用 $\rho(\OmX)=2\rho(\Theta)$ 共轭，我们得到 $B\subset G$ 的一个 $\Theta$-扭形式 $B^{\Theta}$。我们令 $U^{\Theta}\subset B^{\Theta}$ 为其单幂根。对 $G$ 的任一单根 $\alpha$，我们得到相应的投影
\[
\alpha\colon U^{\Theta}\to\OmX.
\]

\srcpage{41}令 $\psi\colon U^{\Theta}\to\OmX$ 为这些映射在全部单根上的和。过渡到分类空间并取从 $X$ 出发的映射叠，这诱导出一个映射
\[
\Bun^{\Theta}_{U}=\Bun_{U^{\Theta}}\to R\Gamma(X,\OmX)[1],
\]
其中我们滥用记号，把 $k$-向量空间与其相应的仿射空间等同。迹映射 $R\Gamma(X,\OmX)[1]\to k$ 由 Serre 对偶给出，它诱导出一个仍记作 $\psi$ 的映射
\[
\psi\colon \Bun^{\Theta}_{U}\to\A^{1}.
\]

在 $\A^{1}$ 上，通常有指数局部系统 $\exp\in D(\A^{1})$——在 D-模的情形，这是解为 $\exp(t)$ 的指数 D-模；在 étale 情形，这是 Artin--Schreier 层。在 Betti 情形它不存在，但可以添入，由此通向野 Betti 层的理论（Scholze，2025），并恢复 Tamarkin（2018）的一个构造。事实上，$\A^{1}$ 上这样的指数局部系统可以添入任何通常的层论，特别地可以添入动机，由此得到指数动机，参见 Fresán 与 Jossen（2025）以及 Gallauer 与 Pepin Lehalleur（2022）。因此，我们将不失一般性地假设我们的层论支持 $\exp$。

最后，利用对应
\[
\begin{tikzcd}
 & \Bun^{\Theta}_{U} \arrow[dl, "p"'] \arrow[dr, "\psi"] & \\
 \Bun_{G} & & \A^{1},
\end{tikzcd}
\]
我们得到 Whittaker 层
\[
\mathcal{W}=\mathcal{W}_{\psi}=p_{!}\psi^{*}\exp\in D(\Bun_{G}).
\]
这就是 $D(\Bun_{G})$ 中我们所期望的基对象。

\begin{remark}
先验地，可以取由单根诱导的映射的另一个线性组合 $\psi'\colon U^{\Theta}\to\OmX$。然而，只要系数非零，$\psi'$ 就会在 $T$ 的作用下与 $\psi$ 共轭，于是相应的 Whittaker 层也相同。归根结底，这种对 $\psi$ 的无关性还意味着：$\mathcal{W}_{\psi}$ 实际上不依赖层论中指数局部系统的选取，因此也可以在 Betti 设定下定义，而不必过渡到野 Betti 层。
\end{remark}

作为例子，取 $G=\mathrm{SL}_2$。此时 $\Bun^{\Theta}_{U}$ 分类扩张
\[
0\to\Theta\to E\to\Theta^{-1}\to 0;
\]
在同构意义下，这些扩张由其在 $\A^{1}$ 中的扩张类分类，这里再次用到 $H^{1}(X,\Theta^{2})=H^{1}(X,\OmX)=k$。每个扩张的自同构群是 $H^{0}(X,\OmX)$。

注意，若 $g\ge 2$，则 $\Theta$ 丰沛，$E$ 被 $\Theta$ 与 $\Theta^{-1}$ 给出的滤过被迫是 Harder--Narasimhan 滤过。事实上，对任意 $G$ 与 $g\ge 2$，映射 $p\colon\Bun^{\Theta}_{U}\to\Bun_{G}$\srcpage{42}分解为 $\Bun^{\Theta}_{U}\to\Bun^{\Theta}_{U}/T\to\Bun_{G}$，其中第二个映射是局部闭浸入，它分类如下轨迹：$G$-丛的 Harder--Narasimhan 滤过成为到 Borel 的约化，且相应的 $T$-丛同构于 $2\rho(\Theta)$。于是，对 $g\ge 2$，Whittaker 层极其显式。这一显式的知识实际上是几何 Langlands 对应证明最后一步的一个关键要素。

我们现在可以定义函子
\[
D_{\qc}(\LocSys_{\Gd})\to D_{\Nilp}(\Bun_G)\colon A\mapsto A\star\mathcal{P}(\mathcal{W}),
\]
其做法是：对 Whittaker 层施加 Beilinson 投影子，再使用谱作用。在 de Rham 情形，省去 Beilinson 谱投影子的施用，而使用整个 $D(\Bun_G)$。

\begin{definition}
粗 Langlands 函子是右伴随
\[
\mathbb{L}^{\coarse}_{G}\colon D_{\Nilp}(\Bun_G)\to D_{\qc}(\LocSys_{\Gd}).
\]
\end{definition}

由于谱作用保持紧对象，该右伴随仍保余极限。真正的 Langlands 函子应当是一个函子
\[
\mathbb{L}_{G}\colon D_{\Nilp}(\Bun_G)\to\IndCoh_{\Nilp}(\LocSys_{\Gd}),
\]
它映入凝聚复形的 Ind-范畴，这些复形在 Arinkin 与 Gaitsgory（2015）意义下的奇异支撑落在幂零锥中。\footnote{这与 $D_{\Nilp}(\Bun_G)$ 形式地类似，但两种奇异支撑的概念不同——分别指可构造层与凝聚层。}有一个遗忘函子
\[
\Psi\colon \IndCoh_{\Nilp}(\LocSys_{\Gd})\to D_{\qc}(\LocSys_{\Gd}),
\]
它是嵌入 $\Coh_{\Nilp}(\LocSys_{\Gd})\subset D_{\qc}(\LocSys_{\Gd})$ 的 Ind-扩张。函子 $\Psi$ 对自然的 $t$-结构是 $t$-正合的，并在最终余连通部分诱导同构——换言之，$\IndCoh_{\Nilp}$ 与 $D_{\qc}$ 之间的差别是 $\infty$-连通的。

\begin{theorem}[Gaitsgory 与 Raskin，2025b]
存在唯一的保余极限函子
\[
\mathbb{L}_{G}\colon D_{\Nilp}(\Bun_G)\to\IndCoh_{\Nilp}(\LocSys_{\Gd}),
\]
使得 $\mathbb{L}^{\coarse}_{G}=\Psi\circ\mathbb{L}_{G}$，并且它把紧对象送到最终余连通的对象。
\end{theorem}

\begin{remark}
让我们简要说明为何扩张是唯一的。由于所有函子都保余极限，它们完全由在紧对象上的行为决定。现在 Gaitsgory 与 Raskin（2025b）证明，$D(\Bun_G)$ 的紧对象在 $\mathbb{L}^{\coarse}_{G}$ 下的像最终余连通。由于 $\Psi$ 在最终余连通对象上是等价，我们就得到所期望的唯一提升。
\end{remark}

\srcpage{43}事实证明，人们可以在某种意义上显式地理解 Hecke 作用——从而谱作用——对 Whittaker 层做了什么。这用到几何 Casselman--Shalika 公式，它是几何 Satake 等价的一个近亲。回忆几何 Satake 等价只是阿贝尔范畴层面的等价；在导出层面，陈述变得更为微妙。反过来，几何 Casselman--Shalika 公式是一个完全的导出等价。

让我们陈述其局部版本。考虑 $LU^{\Theta}$ 在仿射 Grassmann 簇 $\Gr^{\Theta}_{G}$ 上的作用。这里 $\Gr^{\Theta}_{G}$ 是仿射 Grassmann 簇的 $\Theta$-扭形式 $LG^{\Theta}/L^{+}G^{\Theta}$；它参数化形式圆盘上的 $G$-挠子，连同与去心圆盘上的 $G$-丛 $2\rho(\Theta)$ 的一个等同。注意仍有一个映射
\[
\psi\colon LU^{\Theta}\to\A^{1},
\]
其做法是先对到单根空间的投影求和，得到到 $L\OmX$ 的映射，再用留数映射。在 $\A^{1}$ 上，如前有指数局部系统，于是我们考虑
\[
D^{!}(LU^{\Theta}\backslash{}^{\psi}\Gr^{\Theta}_{G})\subset D^{!}(\Gr^{\Theta}_{G}),
\]
即 $\Gr^{\Theta}_{G}$ 上在 $LU^{\Theta}$ 作用下 $\psi$-等变的 $!$-层的满子范畴。把 $\Gr^{\Theta}_{G}$ 写成有限维 Schubert 簇 $\Gr^{\Theta}_{G,\le\mu}$ 的余极限，就有
\[
D^{!}(\Gr^{\Theta}_{G})=\lim_{!}D(\Gr^{\Theta}_{G,\le\mu})=\operatorname{colim}_{!}D(\Gr^{\Theta}_{G,\le\mu}).
\]

\begin{remark}
$!$-层的出现或许可以看作一个迹象，表明人们必须使用 $(-_{!},-_{*})$-形式化。然而，$!$-层在 $(-_{*},-_{!})$-形式化中也有自然的意义。让我们先忘掉 $\Theta$-扭。我们有仿射 Grassmann 簇的自然的 transmutation
\[
[\Gr_{G}]_{D}=\operatorname{colim}_{\mu}[\Gr_{G,\le\mu}]_{D}.
\]
投影 $[\Gr_{G}]_{D}\to[*]_{D}$ 是 $1$-étale 的，但不是 $1$-proper 的。因此，在范畴层的层面，拉回有一个性质良好的左伴随，但没有性质良好的右伴随。所以，从 $[\Gr_{G}]_{D}$ 能构造出的好的层范畴，是把左伴随施用于 $[\Gr_{G}]_{D}$ 上的单位范畴层而得到的，这给出
\[
D^{!}(\Gr_{G})=\operatorname{colim}_{!}D(\Gr_{G,\le\mu}).
\]
类似地，$LU^{\Theta}$ 通过先 transmute 其有限维截断、再取逆向与正向极限，具有自然的 transmutation $[LU^{\Theta}]_{D}$。分类空间 $[*]_{D}/[LU^{\Theta}]_{D}$ 仍是 $1$-prim 的。\footnote{这归约到 $*/[L\A^{1}]_{D}$，在那里可以用 Cartier 对偶的一个应用来证明。}考察复合映射
\[
[LU^{\Theta}]_{D}\backslash[\Gr^{\Theta}_{G}]_{D}\xrightarrow{f}[LU^{\Theta}]_{D}\backslash[*]_{D}\xrightarrow{g}[*]_{D},
\]
我们可以定义
\[
D^{!}(LU^{\Theta}\backslash{}^{\psi}\Gr^{\Theta}_{G})=g_{1*}(f_{1\sharp}\mathbf{1}\otimes\psi),
\]
\srcpage{44}其中 $\psi$ 表示 $[LU^{\Theta}]_{D}\backslash[*]_{D}$ 上的可逆范畴层。它有自然的全忠实遗忘函子
\[
D^{!}(LU^{\Theta}\backslash{}^{\psi}\Gr^{\Theta}_{G})\to D^{!}(\Gr^{\Theta}_{G}),
\]
后者有左伴随。
\end{remark}

要给 $D^{!}(LU^{\Theta}\backslash{}^{\psi}\Gr^{\Theta}_{G})$ 一个干净的定义实际上很微妙，因为 $\Gr^{\Theta}_{G}$ 的归纳有限维本性与 $LU^{\Theta}$ 的归纳（无限维）本性混在一起。上一注的构造表明，Gestalten 的几何能够处理这种复杂性。

几何 Casselman--Shalika 公式给出 $D^{!}(LU^{\Theta}\backslash{}^{\psi}\Gr^{\Theta}_{G})$ 的一个描述。注意，从几何上看，Iwasawa 分解
\[
LG=LB\cdot L^{+}G=\bigsqcup_{\mu\in X_{*}(T)}LU\mu(t)L^{+}G
\]
表明商 $LU^{\Theta}\backslash\Gr^{\Theta}_{G}$ 具有按 $\mu\in X_{*}(T)$ 编号的层片的分层，层片由
\[
LU^{\Theta}\mu(t)L^{+}G^{\Theta}
\]
给出。并非所有层片都承载 $\psi$-等变的层；可以验证，这迫使 $\mu$ 支配。若 $\mu$ 支配，层片上的层等价于 $D(e)$。特别地，$\mu=0$ 给出基对象
\[
\Delta\in D^{!}(LU^{\Theta}\backslash{}^{\psi}\Gr^{\Theta}_{G}).
\]
另一方面，几何 Satake 范畴通过卷积作用在 $D^{!}(LU^{\Theta}\backslash{}^{\psi}\Gr^{\Theta}_{G})$ 上。

\begin{theorem}[几何 Casselman--Shalika]
作用
\[
D_{\qc}(*/\Gd)\to D^{!}(LU^{\Theta}\backslash{}^{\psi}\Gr^{\Theta}_{G})\colon V\mapsto\Sat(V)\star\Delta
\]
是范畴的等价。（事实上，在适当的解释下，得到 $[*]_{D}$ 上范畴层的一个同构。）
\end{theorem}

更一般地，对任何有限集 $I$，我们可以构造带 $LU^{\Theta}$-作用的 $\Gr^{\Theta}_{G,X^{I}}$，并把 $D^{!}(LU^{\Theta}\backslash{}^{\psi}\Gr^{\Theta}_{G,X^{I}})$ 定义为 $[X^{I}]_{D}$ 上的范畴层。在对角线之外，由同样的配方，它同构于 $D_{\qc}(*/\Gd^{I})$ 上的常范畴层；在对角线上，人们得到 $I$ 的商对应的同样情形。把所有函子性都处理妥当，就得到几何 Casselman--Shalika 公式的如下形式。

\srcpage{45}\begin{corollary}
设 $\Whit(G)$ 表示函子
\[
I\mapsto D^{!}(LU^{\Theta}\backslash{}^{\psi}\Gr^{\Theta}_{G,X^{I}})
\]
（视为 $[*]_{D}$ 上的范畴层）的余端。则 $D_{\qc}(*/\Gd)_{\Ran}$ 作用在 $\Whit(G)$ 上，且对 $\Delta$ 的作用诱导出等价
\[
\Whit(G)\cong D_{\qc}(*/\Gd)_{\Ran}.
\]
\end{corollary}

现在，对任何有限集 $I$，令 $\Bun^{U^{\Theta}}_{G,X^{I}}$ 表示 $G$-丛（等价地，经 $2\rho(\Theta)$ 作扭之后，即 $G^{\Theta}$-丛）的模空间，连带一个在 $I$ 点之外到 $U^{\Theta}$ 的约化。有对应
\[
\begin{tikzcd}
 & \Bun^{U^{\Theta}}_{G,X^{I}} \arrow[dl] \arrow[dr] & \\
 LU^{\Theta}\backslash\Gr^{\Theta}_{G,X^{I}} & & \Bun_{G}\times X^{I},
\end{tikzcd}
\]
或在 transmutation 之后类似的图表。此外，由 $\psi$ 给出的 $LU^{\Theta}\backslash*$ 上的范畴层，拉回到 $\Bun^{U^{\Theta}}_{G,X^{I}}$ 后变为平凡，因为在整体曲线上，留数之和为零。于是，由拉推我们得到一个函子\footnote{第二个函子是到 $D^{!}(\Bun_G)=D(\Bun_G)$ 的 $!$-推前。第一个函子用到：$\Bun^{U^{\Theta}}_{G,X^{I}}$ 是 $LU^{\Theta}\backslash\Gr^{\Theta}_{G,X^{I}}$ 沿映射 $*/L^{\mathrm{alg}}U^{\Theta}\to*/LU^{\Theta}$ 的拉回，其中 $L^{\mathrm{alg}}U^{\Theta}$ 把 $R$ 送到 $U^{\Theta}(X_{R}\setminus\bigcup_{i}\Gamma_{x_{i}})$；并且 $L^{\mathrm{alg}}U^{\Theta}$ 是 $0$-suave 的，于是 $*/L^{\mathrm{alg}}U^{\Theta}\to*$ 是 $0$-suave 且对角线 $0$-suave，从而在范畴层的层面，左、右伴随一致。}
\[
D^{!}(LU^{\Theta}\backslash{}^{\psi}\Gr^{\Theta}_{G,X^{I}})\to D^{!}(\Bun^{U^{\Theta}}_{G,X^{I}})\to D(\Bun_{G}).
\]
对全部 $I$ 取余端，我们得到函子
\[
\mathrm{Poinc}_{!}\colon \Whit(G)\to D(\Bun_{G}),
\]
它是 $D_{\qc}(*/\Gd)_{\Ran}$-线性的。由于源等价于 $D_{\qc}(*/\Gd)_{\Ran}$，这个函子等价于函子
\[
D_{\qc}(*/\Gd)_{\Ran}\to D(\Bun_{G})\colon A\mapsto A\star\mathcal{W}_{\psi},
\]
后者由对 Whittaker 层的谱作用给出；但现在它有一个纯粹在自守侧的描述！
